% ===== 1 Introduction =====

This report analyses the paper by \citet{binstok2024} and addresses its research question, methods and results. We focus on the replication of the ALLBUS analysis and its main findings, and we give a short overview of our replication, our critique and our extension of the paper.

\subsection*{Research Problem}

The paper addresses the question of whether personal prejudices shape political views. The existing literature on immigration, race and the War on Terror finds a correlation between prejudice and political views, but provides little causal evidence. A central problem is that prejudice is difficult to measure in surveys: social desirability creates a bias, so that respondents hide their prejudices and their answers do not reflect their true views \citep[pp.~2329--2331]{binstok2024}.

The paper applies this question to a prominent public debate: whether criticism of Israel is motivated by antisemitism. Critics of Israel argue that the accusation of antisemitism is used to silence legitimate criticism, while Jewish advocacy groups argue that parts of this criticism draw on antisemitic stereotypes.
Empirical evidence on the link between antisemitism and views on Israel is limited and not causal \citep[p.~2334]{binstok2024}. \citet[p.~550]{kaplan2006}, for example, show that anti-Israel sentiment predicts antisemitic views, but explicitly make no assumption about causality.

\subsection*{Research Question}

The central research question is therefore whether antisemitic beliefs cause more negative views about Israel \citep[pp.~2331--2332]{binstok2024}. With their methods, the authors aim to identify a systematic link between antisemitism and strong criticism of Israel, but they explicitly do not try to identify the motives of individual persons \citep[p.~2332]{binstok2024}. In this way, the paper contributes causal evidence to a literature on biases and political views that so far mainly documents correlations \citep[pp.~2330--2331]{binstok2024}.

\subsection*{Empirical Approach and Main Findings}

To answer this question, the authors use two independent data sources, each with its own instrumental variable (IV). The first source is Twitter data from the United Kingdom between 2013 and 2022, a panel of 355 districts observed over 3,531 days. For each district and day, they count the negative tweets about Jews and about Israel and compare them.
The authors estimate this relationship with ordinary least squares (OLS) and IV, including district and time fixed effects. As an instrument, they use online chatter about Tottenham Hotspur, a football club that is widely perceived as a ``Jewish'' team but has no connection to Israel \citep[pp.~2333, 2335--2337]{binstok2024}.

The second source, on which this report focuses, is the German General Social Survey (ALLBUS), a cross-section of individuals from 2016.
Respondents rate two statements on a scale from 1 (completely disagree) to 7 (completely agree): ``Jews have too much influence in the world'', which measures antisemitism, and ``It is unjust that Israel is taking away land from the Palestinians'', which measures criticism of Israel. The authors test whether people who agree more with the first statement also agree more with the second.

The authors again use OLS and IV. The instrument is being born between 1930 and 1945, i.e.\ exposure to Nazi indoctrination during childhood, following the idea of \citet{voigtlaender2015}, compared with respondents born between 1946 and 1960 \citep[pp.~2345--2346, 2355--2358]{binstok2024}. As a robustness check, they use a 2004 survey of the Anti-Defamation League (ADL) covering ten European countries, again with birth cohorts as instruments \citep[pp.~2365--2366]{binstok2024}.

In both analyses, the IV estimates are considerably larger than the OLS estimates: 
On Twitter, the OLS coefficient is 0.30 and the IV coefficient 0.75; in the ALLBUS data, the OLS coefficient is about 0.16, while the main IV estimate is 0.85 \citep[pp.~2338, 2340, 2353, 2359]{binstok2024}. The authors explain this gap with under-reporting: respondents conceal their antisemitic views in surveys but express them through their views on Israel \citep[p.~2364]{binstok2024}. At the same time, they stress that about half of the respondents with the most critical views of Israel show no sign of antisemitism \citep[p.~2351]{binstok2024}.

\subsection*{Our Replication and Extension}

Our replication focuses on the ALLBUS analysis. We reproduce the means of all survey statements in Table~6 (column~1), the OLS results in Table~7, the first stage and the reduced form in Table~9 (third panel), the IV results in Table~10 and Figure~7 of the paper; the reasons for this selection are set out in Section~\ref{sec:replication}. All numbers match the published results exactly.

While replicating, we noticed that the first stage is weak: in the main specification, the $F$-statistic is only 10.5, which the paper does not report. Such a low $F$-statistic means that the instrument has only a small effect on antisemitic views, so that the IV estimate can be biased. According to the rule of thumb of \citet{lee2022}, the $F$-statistic should exceed about 104.7 for the usual critical value of $\pm 1.96$ to be valid. In our extension, we therefore use the tF procedure, which remains valid even when the instrument is weak. With tF, the main estimate is no longer significant at the 5\% level ($p = 0.14$). The Anderson--Rubin confidence interval \citep{anderson1949}, which is also robust to weak instruments, still excludes zero, but it is very wide (0.33 to 2.18). Moreover, a direct effect of the birth cohort on views about Israel of only 0.13 points on the 1--7 scale would be enough to make the estimate insignificant \citep{conley2012}.



% ---------------------------------------------------------------------------
% RAW MATERIAL: draft of our team member (v2). Not shown in the PDF.
% Use as a source only; rewrite in your own words and shorten.
% ---------------------------------------------------------------------------
% \par\textbf{Introduction}\par
%
% This chapter summarises the paper by Binstok, Gould and Kaplan (2024) along the four questions of Problem Set 3, part (a): the problem the authors address, their research question, their empirical approach and their main findings. Each statement is referenced to the page and line of the published article.
%
% \par\textbf{Research Problem}\par
%
% \begin{itemize}
% \item
%   The paper addresses the question of whether personal prejudices shape political views, in particular views on policies that predominantly affect a minority group.\footnote{(Binstok, Gould, and Kaplan 2024, pp. 2329--2330, l. 23--24 and l. 1--9)}
% \item
%   The existing literature on race, immigration, welfare and the War on Terror documents correlations between prejudice and policy views, but provides little causal evidence.\footnote{(Binstok, Gould, and Kaplan 2024, pp. 2330--2331, l. 11--22 and l. 1--4)}
%
%   \begin{itemize}
%   \item
%     Two obstacles are identified: prejudice is difficult to measure in surveys, and social desirability bias leads respondents to conceal their true views.\footnote{(Binstok, Gould, and Kaplan 2024, p. 2330, l. 15--20)}
%   \end{itemize}
% \item
%   The specific context is the public debate on whether criticism of Israel is driven by antisemitism.\footnote{(Binstok, Gould, and Kaplan 2024, pp. 2331--2332, l. 8--18 and l. 1--17)}
%
%   \begin{itemize}
%   \item
%     Critics of Israel argue that accusations of antisemitism are used to silence legitimate criticism.\footnote{(Binstok, Gould, and Kaplan 2024, p. 2331, l. 10--13; p. 2332, l. 14--17)}
%   \item
%     Jewish advocacy groups and the IHRA definition hold that parts of anti-Israel rhetoric draw on antisemitic tropes.\footnote{(Binstok, Gould, and Kaplan 2024, p. 2331, l. 13--18; p. 2332, l. 1--13)}
%   \end{itemize}
% \item
%   Empirical evidence on the link between antisemitism and views on Israel is scarce and not causal; Kaplan and Small (2006), for instance, document a correlation only.\footnote{(Binstok, Gould, and Kaplan 2024, p. 2334, l. 18--21 and fn. 9, l. 42--46)}
% \end{itemize}
%
% \par\textbf{Research Question}\par
%
% \begin{itemize}
% \item
%   The central research question is whether antisemitic beliefs cause more negative views about Israel.\footnote{(Binstok, Gould, and Kaplan 2024, p. 2330, l. 10--11; p. 2331, l. 4--5)}
% \item
%   The authors aim to identify a systematic, aggregate link between antisemitism and strong criticism of Israel; they explicitly do not assess the motives of any individual critic.\footnote{(Binstok, Gould, and Kaplan 2024, p. 2332, l. 18--25)}
% \item
%   The question is embedded in the broader literature on how biases affect views on public policy.\footnote{(Binstok, Gould, and Kaplan 2024, pp. 2370--2371, l. 39--40 and l. 1--9)}
% \end{itemize}
%
% \par\textbf{Empirical Approach}\par
%
% \begin{itemize}
% \item
%   The paper combines two independent empirical strategies, each relying on its own instrumental variable.\footnote{(Binstok, Gould, and Kaplan 2024, p. 2332, l. 26--28; p. 2335, l. 11--14)}
% \item
%   \textbf{Strategy 1, Twitter (United Kingdom, 2013--2022):} a panel of 355 districts observed over 3,531 days, estimated by OLS and IV with district and time fixed effects.\footnote{(Binstok, Gould, and Kaplan 2024, p. 2333, l. 6--20; p. 2335, l. 16--25)}
%
%   \begin{itemize}
%   \item
%     Instrument: online chatter about Tottenham Hotspur, a football club widely perceived as a ``Jewish'' team but unrelated to Israel.\footnote{(Binstok, Gould, and Kaplan 2024, pp. 2332--2333, l. 29--35 and l. 1--8)}
%   \end{itemize}
% \item
%   \textbf{Strategy 2, German General Social Survey (ALLBUS 2016):} a cross-section of individuals, estimated by OLS with successively larger sets of controls and subsequently by IV.\footnote{(Binstok, Gould, and Kaplan 2024, pp. 2333--2334, l. 21--34 and l. 1--2)}
%
%   \begin{itemize}
%   \item
%     Instrument: being born between 1930 and 1945, i.e. exposure to Nazi indoctrination during childhood (following Voigtländer and Voth 2015), compared with respondents born 1946--1960.\footnote{(Binstok, Gould, and Kaplan 2024, p. 2334, l. 3--9)}
%   \end{itemize}
% \item
%   Robustness: a 2004 survey of the Anti-Defamation League (ADL) covering ten European countries, again with birth cohort as instrument.\footnote{(Binstok, Gould, and Kaplan 2024, p. 2334, l. 12--13; p. 2365, l. 39--41)}
% \end{itemize}
%
% \par\textbf{Main Findings}\par
%
% \begin{itemize}
% \item
%   \textbf{Twitter:} negative tweets about Jews are strongly correlated with negative tweets about Israel at the district-day level.\footnote{(Binstok, Gould, and Kaplan 2024, p. 2333, l. 12--13)}
%
%   \begin{itemize}
%   \item
%     The OLS coefficient is 0.297, the IV coefficient 0.748, i.e. more than twice as large, with a first-stage F-statistic above 25.\footnote{(Binstok, Gould, and Kaplan 2024, p. 2340, l. 7--12)}
%   \item
%     In a placebo test, negative tweets about Jews have no effect on negative tweets about China.\footnote{(Binstok, Gould, and Kaplan 2024, p. 2333, l. 18--20; p. 2340, l. 35--38)}
%   \item
%     The effect on positive tweets about Israel is small.\footnote{(Binstok, Gould, and Kaplan 2024, p. 2340, l. 30--34)}
%   \end{itemize}
% \item
%   \textbf{ALLBUS, OLS:} agreement that ``Jews have too much influence'' is positively associated with agreement that ``Israel is unjust to take Palestinian land''; the coefficient rises from 0.118 to 0.162 as controls are added (Table 7, upper panel, p. 2353).\footnote{(Binstok, Gould, and Kaplan 2024, p. 2352, l. 11--15 and l. 34--37)}
% \item
%   \textbf{ALLBUS, IV:} the estimated effects are considerably larger than the OLS estimates and also increase with the set of controls.\footnote{(Binstok, Gould, and Kaplan 2024, p. 2334, l. 9--12)}
%
%   \begin{itemize}
%   \item
%     In the main specification, the coefficient is 0.850 with a standard error of 0.339 (Table 10, column 4, p. 2359; repeated in Table 11).\footnote{(Binstok, Gould, and Kaplan 2024, p. 2362, l. 12--13)}
%   \item
%     In terms of magnitude, an increase of one standard deviation in antisemitism (factor measure) raises anti-Israel views by about 1.10 standard deviations under IV, compared with 0.18 under OLS.\footnote{(Binstok, Gould, and Kaplan 2024, p. 2364, l. 18--27)}
%   \end{itemize}
% \item
%   The authors attribute the gap between IV and OLS to under-reporting: respondents conceal antisemitic views (measurement error) but express them through their views on Israel.\footnote{(Binstok, Gould, and Kaplan 2024, p. 2364, l. 27--35; p. 2370, l. 34--38)}
% \item
%   The authors stress an important qualification: about half of the respondents with the most negative views of Israel show no sign of antisemitism.\footnote{(Binstok, Gould, and Kaplan 2024, p. 2351, l. 8--20; p. 2365, l. 35--38)}
% \item
%   In a counterfactual exercise, capping antisemitism at level 3 lowers the mean anti-Israel view by 0.76, i.e. by roughly half a standard deviation.\footnote{(Binstok, Gould, and Kaplan 2024, p. 2365, l. 12--19)}
% \item
%   The ADL data yield a similar pattern, with IV estimates exceeding OLS estimates and rising with the set of controls.\footnote{(Binstok, Gould, and Kaplan 2024, p. 2370, l. 23--29)}
% \end{itemize}
%
% \par\textbf{Replication and Extension}\par
%
% \begin{itemize}
% \item
%   \textbf{Replication:} the analysis focuses on the ALLBUS results, namely Table 6 (column 1), Table 7, the first stage in Table 9 (third panel, column 1), Table 10 and Figure 7. The reasons for this selection are set out in Chapter 3.
% \item
%   \textbf{Status:} all selected estimates have been reproduced exactly in R from the raw ALLBUS file, including the main estimate of 0.850 (0.339).
% \item
%   \textbf{Extension:} still to be specified. A natural candidate is the strength of the first stage, which falls to about 10 in the main specification (Section 3.4).
% \end{itemize}
